% ============================================================================
\section{Diseño de la investigación}

\subsection{Tipo y diseño}

El estudio es cuantitativo, aplicado y no experimental. Combina el análisis histórico de las sucesivas versiones publicadas del PBI peruano con simulaciones numéricas controladas para poner a prueba los modelos en escenarios conocidos. Todas las predicciones se evaluaron fuera de muestra y respetando estrictamente el paso del tiempo. Se fijó por adelantado una arquitectura cuántica principal para evitar elegir el mejor resultado a conveniencia después de ver los números.

\subsection{Material, métodos y técnicas}

\subsubsection{Material: población y muestra}

La población abarca todas las publicaciones mensuales del PBI del Perú registradas de manera verificable en las Notas Semanales del BCRP. Tras la limpieza inicial se reunieron \(44\,307\) datos por actividad económica correspondientes a nueve sectores a lo largo de \(384\) boletines mensuales continuos (enero de 1994 a diciembre de 2025), cubriendo \(406\) meses de actividad (enero de 1992 a octubre de 2025).

Para el PBI agregado, al medir la revisión a un plazo fijo de once meses (\(H=11\)) y aplicar cinco reglas de calidad quedaron \(350\) observaciones limpias (\(374\) antes de retirar los cambios de base) y \(331\) observaciones al plazo de doce meses:
\begin{enumerate}
	\item identificar sin ambigüedad a qué mes corresponde cada cifra;
	\item exigir una primera publicación oportuna (retirando los \(21\) meses de 1992--1993 con retrasos iniciales largos);
	\item contar con el boletín correspondiente \(H\) meses después;
	\item excluir los \(24\) meses cruzados por los cambios de año base de julio de 2000 y marzo de 2014;
	\item comprobar la ausencia de filas repetidas y códigos especiales sin dato válido.
\end{enumerate}
El detalle completo de filtros y meses descartados figura en el Anexo~\ref{app:auditoria}.

\subsubsection{Fuentes y disponibilidad de información}

Se trabajó sobre la base histórica abierta \emph{Peru GDP Real-Time Dataset (1994--2025)} de \textcite{cruz2025}. Al inspeccionar los boletines semanales (con \(384\) meses únicos luego de quitar \(150\) reportes repetidos), se comprobó que las tablas indican el mes y el número de nota semanal, pero no el día exacto de salida. Por ello, los cálculos trabajan bajo una aproximación mensual en tiempo real.

\subsubsection{Métodos de análisis y pruebas estadísticas}

Para responder cada pregunta del estudio se combinaron cinco grupos de herramientas:
\begin{itemize}
	\item medidas descriptivas básicas (promedio, mediana, dispersión y valores extremos);
	\item pruebas estadísticas clásicas para detectar sesgo, memoria en el tiempo \parencite{ljung1978}, eficiencia informativa \parencite{mincer1969}, relaciones curvas \parencite{ramsey1969} y cambios de volatilidad entre épocas \parencite{engle1982}, ajustando los errores estándar para tomar en cuenta la cercanía temporal entre meses \parencite{newey1987};
	\item referentes clásicos de pronóstico: modelos autorregresivos, regresión Ridge \parencite{hoerl1970}, factores sectoriales \parencite{stock2002} y árboles de decisión acotados \parencite{friedman2001,chen2016};
	\item reservorios y transformaciones no lineales de tamaño equivalente: el modelo cuántico QRC \parencite{fujii2017}, la red clásica ESN \parencite{jaeger2001} y las características aleatorias RFF \parencite{rahimi2007};
	\item pruebas modernas para comparar pronósticos sin dejarse engañar por el azar al comparar muchos modelos a la vez: contraste de errores \parencite{diebold1995,harvey1997,clark2007}, remuestreo por bloques temporales \parencite{politis1992}, ajustes por comparaciones múltiples \parencite{holm1979,benjamini1995} y selección del grupo de mejores modelos empatados \parencite{hansen2011}.
\end{itemize}

\subsubsection{Recolección y cuidado de los datos}

Los archivos originales del repositorio se mantuvieron intactos en todo momento. Se automatizó la lectura, ordenamiento cronológico y verificación matemática mediante códigos reproducibles y huellas digitales de control (sumas criptográficas SHA-256).

\subsection{Estrategia paso a paso del estudio}

\subsubsection{Limpieza del archivo histórico}

Se ordenó la base en una matriz donde cada fila representa un mes medido y cada columna una fecha de publicación, confirmando que ninguna celda mostraba información adelantada al futuro. Además, se separaron las revisiones estadísticas cotidianas de los saltos causados por los cambios de año base.

\subsubsection{Elección del horizonte y diagnóstico inicial}

El plazo principal de once meses (\(H=11\)) se fijó exclusivamente por continuidad de datos antes de probar ningún modelo predictivo, dejando \(H=12\), \(6\) y \(3\) meses como respaldo. Sobre la serie depurada se revisó si las revisiones tienen sesgo sistemático, ciclos repetitivos por mes del año o cambios bruscos de comportamiento antes y después de la pandemia.

\subsubsection{Modelos clásicos de referencia}

Se estimaron ocho referentes clásicos usando únicamente la información disponible en cada mes evaluado: referencia nula (predecir cero), promedio acumulado histórico, promedio móvil de sesenta meses, modelos autorregresivos (AR y ARX), regresión Ridge sobre veinte variables observables, factores dinámicos de seis sectores continuos y árboles de decisión acotados implementados en Python ante la ausencia de librerías externas.

\subsubsection{Modelo cuántico y controles clásicos equivalentes}

El reservorio cuántico se implementó mediante simulación exacta en Python de cinco qubits acoplados (sin ruido de laboratorio), midiendo veinte señales cuánticas simples que alimentan una regresión Ridge final. Para comparar en justicia con herramientas clásicas de exactamente el mismo tamaño de lectura (\(m=20\)), se evaluaron una red clásica ESN de veinte neuronas, una variante ESN ampliada de sesenta y cuatro neuronas y el mapeo estático de Fourier (RFF), repitiendo cada algoritmo cinco veces con distintas semillas aleatorias para comprobar su estabilidad \parencite{schuld2021}.

\subsubsection{Simulación realista mes a mes sin ver el futuro}

Los doce modelos compitieron mes a mes entre febrero de 2003 y noviembre de 2024 (\(250\) fechas en total). En cada fecha \(o\), cada modelo entrenaba solo con meses pasados cuya revisión ya había terminado:
\[
\widehat{Y}_{o+1|o}=f_o(\mathcal{I}_o).
\]
Los últimos cinco años (\(60\) meses, de diciembre de 2019 a noviembre de 2024) quedaron guardados aparte hasta el final, de modo que ningún modelo pudiera ajustarse mirando la pandemia de antemano.

\subsubsection{Escenarios simulados y potencia estadística}

Como en la historia real solo tenemos \(250\) meses de prueba, también se crearon doce escenarios artificiales calibrados a partir de los datos peruanos mediante la ecuación:
\begin{equation}
	r_{t,j}=\mu_j+\phi_j r_{t-1,j}
	+\gamma_j s_t+\delta_j B_t+\varepsilon_{t,j},
	\label{eq:dgp}
\end{equation}
tras \(300\) remuestreos de incertidumbre por bloques \parencite{politis1992}. Variando el largo de la serie, el grado de memoria, la fuerza de las señales y la presencia de choques fuertes (\(20\) simulaciones por escenario), se midió cuándo logra ganar QRC y cuánta diferencia real en precisión se necesita en el Perú para detectarla con seguridad (Anexo~\ref{app:potencia}).

\subsubsection{Regla de parada para el análisis económico}

La valoración económica mediante una regla cuadrática sencilla de costo por errores macroeconómicos:
\begin{equation}
	L_t=(\pi_t-\pi^\star)^2+\lambda_y \widehat{gap}_t^2,
	\label{eq:loss}
\end{equation}
se programó bajo una regla estricta: solo se calcularía si QRC vencía de forma clara y robusta a los modelos clásicos equivalentes. Como esa ventaja no apareció, el análisis económico no se ejecutó para evitar inventar beneficios irreales.

\subsubsection{Cómo se midió y comparó el error de los pronósticos}

Para medir qué tan cerca estuvieron los pronósticos de las revisiones reales se usaron principalmente la raíz del error cuadrático medio (RMSE, que castiga más los errores grandes) y el error absoluto medio (MAE, el error promedio en puntos porcentuales):
\[
\RMSE=\sqrt{\frac{1}{T_o}\sum_{t=1}^{T_o}e_t^2},
\qquad
\MAE=\frac{1}{T_o}\sum_{t=1}^{T_o}|e_t|,
\]
junto con el sesgo promedio y el porcentaje de mejora respecto a una referencia (\(R^2_{\text{OOS}}\)). Se evitó usar errores en porcentaje (MAPE) porque se indeterminan cuando una revisión real vale cero.

\subsection{Reglas para evitar resultados engañosos}

Durante toda la investigación se cuidaron siete principios básicos:
\begin{enumerate}
	\item \textbf{Cero información futura:} ninguna variable ni ajuste usó datos no publicados en la fecha evaluada.
	\item \textbf{Sin elegir después de ver:} el modelo cuántico y sus competidores se definieron antes de abrir los últimos \(60\) meses.
	\item \textbf{Control del sobreajuste:} se limitaron las lecturas a veinte señales y todo parámetro se ajustó dentro del pasado.
	\item \textbf{Comparación justa:} los controles clásicos tuvieron exactamente el mismo tamaño de lectura que QRC.
	\item \textbf{Respeto al orden del tiempo:} las pruebas tomaron en cuenta que meses seguidos comparten información.
	\item \textbf{Separación de cambios de año base:} se aislaron los empalmes metodológicos de 2000 y 2014.
	\item \textbf{Reproducibilidad total:} todo el código, semillas, parámetros y bases originales quedaron guardados sin alteración.
\end{enumerate}

\subsection{Criterios fijados de antemano para decidir}

Como umbral mínimo de relevancia práctica se fijó una reducción relativa de al menos \(2{,}0\%\) en el error RMSE. De acuerdo con los resultados conjuntos, las conclusiones se clasificaron en cuatro categorías claras: \textbf{evidencia favorable} (mejora relevante confirmada tras corregir comparaciones múltiples), \textbf{evidencia desfavorable} (sin mejora o superado por un competidor clásico), \textbf{resultado inconcluso} (margen estadístico amplio que incluye el empate) o \textbf{problema de datos} (si fallara la reconstrucción histórica). La falta de diferencia significativa se interpreta con prudencia como falta de evidencia suficiente, no como prueba automática de igualdad exacta.
